% ============================================================
% Chapter: Concept Selection and Life-Cycle Evaluation
% Assessed MLOs: MLO-1, MLO-2
% Excellent-grade rubric criteria:
% - Concept design selection and evaluation (10%): alternatives are physically
%   credible, screened with explicit criteria, and the chosen concept is justified
%   through engineering evidence rather than preference.
% - Robust evaluation of life-cycle alternatives (5%): temporary stages,
%   installation, maintenance, and end-of-life considerations materially inform
%   the concept choice.
% ============================================================

\section{Concept Selection and Life-Cycle Evaluation}
\pagecheck{8--9}

% Rubric Checklist: Concept design selection and evaluation (10%); Robust evaluation of life-cycle alternatives (5%).
\instructionplaceholder{Concept screening}{Present the concept alternatives that were considered seriously enough to compare. For each alternative, state the structural principle, the primary load path, the main advantages, and the primary concerns or elimination criteria. Avoid vague descriptions; include sketches or free-body diagrams.}

\subsection{Selection criteria and decision logic}
\instructionplaceholder{Required evidence}{Define the criteria used to compare concepts. These should cover technical feasibility, stability or equilibrium, constructability, maintainability, robustness, and project-specific constraints. Explain the weighting or ranking logic and show why the selected concept performs best against the stated criteria.}

\subsection{First-order equilibrium and feasibility}
\instructionplaceholder{Required evidence}{Provide the first-order calculations that establish vertical and horizontal load balance. Show buoyancy, self-weight, mooring reactions, and other dominant actions in a transparent way so the reader can judge whether the concept is physically viable before numerical refinement.}

\subsection{Life-cycle evaluation}
\instructionplaceholder{Required evidence}{Assess the consequences of construction, transport, installation, inspection, maintenance, repair, and decommissioning for the chosen concept. State which temporary stages are critical, what risks or bottlenecks arise during the life cycle, and how these findings strengthen or weaken the selected concept.}
